\chapter{Certifying the Riemann Hypothesis up to bounded height}
\label{cha:rh_up_to}

Finding certified counts of zeros of the zeta function on the critical line 
yields a lower bound of the form $n + 1 \leq N(T)$.
On the other hand, Turing's method can be used to derive an upper bound 
of the form $N(g_n)\leq n + 1$. 
If $T \leq g_n$ then this forces $N(T) = n + 1$ and therefore forces 
all the zeros of the zeta function on the domain 
$\{z | \text{Im}(z) \in (0, T)\}$ to be on the critical line.

\begin{theorem}
\label{thm:rh_up_to}
\lean{riemannHypothesisUpTo}
\leanok
Suppose that $g_n$ is a Gram point satisfying the conditions of 
Theorem~\ref{thm:n_upper_bound}. Suppose that $Z$ is a certified 
list of zeros of the zeta function on the subset $\{1/2 + it\mid 0 < t < g_n\}$ 
of the critical line using a search algorithm. 
Then all non-trivial zeros of the zeta function in the region 
$|\text{Im}(s)| < g_n$ are on the critical line.
\end{theorem}

\begin{proof}
\leanok
\uses{thm:n_upper_bound}
On the real line, the only zeros of $\zeta$ are the trivial zeros.
Moreover, because $\zeta(\overline{s}) = \overline{\zeta(s)}$, it suffices to 
consider the zeros in $\{s \mid 0 \leq \text{Im}(s) < g_n\}$.
The number of such zeros, with multiplicity, is by definition $N(g_n)$.

By Theorem~\ref{thm:n_upper_bound}, $N(g_n)\leq n + 1$. 
By the certified count there are at least $n + 1$ zeros on 
$\{1/2 + it \mid 0 \leq t < g_n\}$.
Therefore all zeros up to height $g_n$ lie on the critical line.
\end{proof}
